Anomaly detection and root cause analysis are crucial in nowadays enterprise IT applications \cite{16_Liu2020_TraceAnomaly,29_Wang2018_CloudRanger}.
Techniques for solving both problems have already been discussed in existing surveys, which however differ from ours as they consider either anomaly detection or root cause analysis.

%% ANOMALY DETECTION - SURVEYS 
Chandola \etal \cite{Chandola2009_SurveyAnomalyDetection} and Akoglu \etal \cite{Akoglu2015_GraphBasedAnomalyDetection} survey the research on detecting anomalies in data.
% In both surveys, the focus is on presenting techniques for detecting significant changes and outliers in data, independently of the actual application for such techniques.
% We differ from both surveys as we fix the application domain: we indeed focus on determining anomalies that can be symptoms of failures of the services forming an application.
% We hence reduce the focus to those anomaly detection techniques that have been designed to process the time series data corresponding to logs, traces, or KPIs monitored on application services.
% This enables us to provide a more detailed description for such techniques, as well as to compare and discuss them under dimensions peculiar to our application domain, \eg whether they can detect application- or service-level anomalies, the cost for their setup, or their accuracy in dynamically changing applications.
\upd{In both surveys, the focus is on the generic problem of significant changes/outlier detection in data, independently of whether this is done online or offline, or what the data is about.
We differ from both surveys as we focus on online anomaly detection in a specific type of data, \viz logs, traces, or KPIs monitored on application services.
This allows us to provide a more detailed description of how such techniques can detect anomalies in running multi-service applications, as well as to compare and discuss them under dimensions peculiar to such problem, \eg whether they can detect application- or service-level anomalies, the cost for their setup, or their accuracy in dynamically changing applications.}
In addition, the studies surveyed by Chandola \etal \cite{Chandola2009_SurveyAnomalyDetection} and Akoglu \etal \cite{Akoglu2015_GraphBasedAnomalyDetection} are dated 2015 at most, hence meaning that various anomaly detection techniques for multi-service applications are not covered by such surveys.
Such techniques have indeed been proposed after 2014, when microservices were first proposed \cite{Lewis2014_Microservices} and gave rise to the widespread use of multi-service architectures \cite{Soldani2018_MicroservicesPainsGains}.
Finally, differently from Chandola \etal \cite{Chandola2009_SurveyAnomalyDetection} and Akoglu \etal \cite{Akoglu2015_GraphBasedAnomalyDetection}, we also survey the analysis techniques that can be used in a pipeline with the surveyed anomaly detection techniques to determine the possible root causes for detected anomalies.

% ROOT CAUSE ANALYSIS - SURVEYS 
Steinder and Sethi \cite{Steinder2004_FaultLocalizationTechniquesComputerNetworks}, Wong \etal \cite{Wong2016_SurveySoftwareFaultLocalization}, and Sole \etal \cite{Sole2017_SurveyRootCauseAnalysis} survey solutions for determining the possible root causes for failures in software systems.
In particular, Steinder and Sethi \cite{Steinder2004_FaultLocalizationTechniquesComputerNetworks} survey techniques for determining the root causes of anomalies in computer networks.
Wong \etal \cite{Wong2016_SurveySoftwareFaultLocalization} instead provide an overview of the techniques allowing to identify the root causes of failures in the source code of a single software program.  
\upd{Both Steinder and Sethi \cite{Steinder2004_FaultLocalizationTechniquesComputerNetworks} and Wong \etal \cite{Wong2016_SurveySoftwareFaultLocalization} hence
differ from our survey as they focus on software systems different than multi-service applications}.
%
In this perspective, the survey by Sole \etal \cite{Sole2017_SurveyRootCauseAnalysis} is closer to ours, since they focus on techniques allowing to determine the root causes of anomalies observed in multi-component software systems, therein included multi-service applications.
The focus of Sole \etal \cite{Sole2017_SurveyRootCauseAnalysis} is however complementary to ours: they analyse the performance and scalability of the surveyed root cause analysis techniques, whilst we focus on more qualitative aspects, \eg what kind of anomalies can be analysed, which analysis methods are enacted, which instrumentation they require, or which actual root causes are determined.
In addition, we also survey existing techniques for detecting anomalies in multi-service application, with the aim of supporting application operators in setting up their pipeline to detect anomalies in their applications and to determine the root causes of observed anomalies.

% ROOT CAUSE ANALYSIS - EVALUATION
Another noteworthy secondary study is that by Arya \etal \cite{Arya2021_EvaluationCausalInferenceAIOps}, who evaluate several state-of-the-art techniques for root cause analysis, based on logs obtained from a publicly available benchmark multi-service applications. 
In particular, they consider Granger causality-based techniques, by providing a first quantitative evaluation of their performance and accuracy on a common dataset of application logs.
They hence complement the qualitative comparison in our survey with a first result along a research direction emerging from our survey, \viz quantitatively comparing the performance and accuracy of the existing root cause analysis techniques. 

% SUMMARY
In summary, to the best of our knowledge, ours is the first survey on the techniques for detecting anomalies that can be symptoms of failures in multi-service applications, while at the same time presenting the analysis techniques for determining possible root causes for such anomalies.
It is also the first survey providing a more detailed discussion of such techniques in the perspective of their use with multi-service applications, \eg whether they can detect and analyse anomalies at application- or service-level, which instrumentation they require, and how they behave when the services forming an application and their runtime environment change over time.